% !TeX root = ../../tesis.tex

\subsection{Salidas estructurales: elección del formato de intercambio}
\label{subsec:estructura-geojson}

Los datos geoespaciales de la red de transporte metropolitana pueden
distribuirse en tres formatos vectoriales de uso extendido en el ámbito urbano:
el \textit{Shapefile} (\textsc{shp}), el \textit{GeoPackage} y el
\textit{GeoJSON}. Cada uno responde a un contexto de uso distinto
\citep{longley2015}, como se sintetiza en la Tabla~
\ref{tab:formatos-geoespaciales}.

\begin{table}[H]
    \centering
    \small
    \begin{tabular}{|p{2.2cm}|p{3.2cm}|p{3.8cm}|p{4.2cm}|}
        \hline
        \textbf{Formato} & \textbf{Tipo de archivo} &
        \textbf{Uso habitual} & \textbf{Limitaciones} \\ \hline
        %
        \textsc{shp} (Shapefile) &
        Múltiples archivos binarios (.shp, .dbf, .shx)
        \citep{esri1998shapefile} &
        Escritorio \textsc{sig}; herencia institucional
        (\textsc{inegi}, \textsc{adip}) &
        Requiere al menos 3 archivos; límite en nombres de campo;
        problemas de codificación en entornos web \\ \hline
        %
        GeoPackage (.gpkg) &
        Contenedor \textsc{sqlite} binario, archivo único
        \citep{ogc2014geopackage} &
        Formato nativo de \textsc{qgis}; análisis local en
        escritorio &
        Binario: no legible como texto; difícil de gestionar con
        control de versiones \\ \hline
        %
        GeoJSON &
        Texto plano \textsc{json}, archivo único
        \citep{butler2016geojson} &
        Servicios web, \api{} REST, visores interactivos &
        Mayor peso en disco para geometrías con millones de
        vértices \\ \hline
    \end{tabular}
    \caption[Comparativa de formatos de datos geoespaciales]{Comparativa
    entre los tres formatos vectoriales más comunes en flujos de trabajo
    urbano.}
    \label{tab:formatos-geoespaciales}
\end{table}

El \textit{Shapefile} continúa presente en entregas de datos institucionales en
México —archivos de \textsc{inegi}, \textsc{adip} y \textsc{semovi} llegan con
frecuencia en este formato—, por lo que su lectura sigue siendo parte del
repertorio cotidiano del urbanista. Sin embargo, para un servicio de datos en
red como \texttt{Apimetro}, el \textit{GeoJSON} ofrece dos ventajas
estructurales sobre sus alternativas: es texto plano —cualquier corrupción se
manifiesta como un error de parseo inmediato, mientras que un archivo binario
dañado puede aparentar ser válido hasta que se intenta abrir en el \textsc{sig}—
y es consumible directamente por herramientas web y entornos de programación sin
conversión previa \citep{butler2016geojson}. Para el urbanista que trabaja en
escritorio, la conversión de \textit{GeoJSON} a \textit{GeoPackage} es una
operación de una sola línea en Python con la biblioteca \textit{GeoPandas}
\citep{tycgis2025}, lo que preserva la interoperabilidad con \textsc{qgis} sin
sacrificar las ventajas del formato de intercambio. Las métricas operativas
necesarias para el motor \textit{VFTModel} —velocidad promedio, frecuencia y
capacidad del vehículo— viajan integradas en las \textit{properties} de cada
objeto GeoJSON, eliminando el procesamiento adicional de cruce entre archivos
\citep{cabrera2026VFTModel}.

\subsection{JSON Descriptivo: la ficha técnica de la red}
\label{subsec:estructura-geojson-descriptivo}

\begin{figure}[H]
    \centering
    % Diagrama: Dos tipos de salida de Apimetro
% Columna izquierda: JSON Descriptivo (atributos alfanuméricos)
% Columna derecha: GeoJSON Geográfico (geometría + métricas)
% Usado en: 4-Capas-Codigo/4-1-Apimetro/4-1-4-Salidas.tex (fig:apimetro-salidas-tipos)
\begin{tikzpicture}[
  cabecera/.style={
    rectangle, draw=unamAzul, fill=unamAzul, text=white,
    minimum width=5.4cm, minimum height=0.9cm,
    align=center, font=\small\bfseries
  },
  cuerpo/.style={
    rectangle, draw=unamAzul!40, fill=unamAzul!8,
    minimum width=5.4cm, text width=5.0cm,
    align=left, inner sep=8pt, font=\footnotesize
  },
  uso/.style={
    rectangle, draw=unamOro!70, fill=unamOro!18,
    minimum width=5.4cm, minimum height=0.75cm,
    align=center, font=\footnotesize\itshape, text=unamAzul
  }
]

%% ── Columna izquierda: JSON Descriptivo ──────────────────────────
\node[cabecera] (jh) at (0, 0) {JSON Descriptivo};

\node[cuerpo, anchor=north] (jb) at (jh.south) {%
  \textbf{¿Qué contiene?}\\[4pt]
  \texttt{linea\_id} · \texttt{nombre}\\
  \texttt{sistema} · \texttt{color\_esp}\\
  \texttt{clasificacion}\\[4pt]
  \textit{Sin geometría}
};

\node[uso, anchor=north] (ju) at (jb.south)
  {tablas · fichas · simbología};

%% ── Columna derecha: GeoJSON Geográfico ──────────────────────────
\node[cabecera] (gh) at (6.6, 0) {GeoJSON Geográfico};

\node[cuerpo, anchor=north] (gb) at (gh.south) {%
  \textbf{¿Qué contiene?}\\[4pt]
  \texttt{coordinates} (Point / MultiLineString)\\
  \texttt{properties}\\
  \quad velocidad · frecuencia · capacidad\\[4pt]
  \textit{Listo para web, sin conversión}
};

\node[uso, anchor=north] (gu) at (gb.south)
  {mapas · ruteo · \textit{VFTModel}};

\end{tikzpicture}

    \caption[Tipos de salida de Apimetro]{Los dos tipos de respuesta de
    la \api{}: JSON descriptivo (atributos alfanuméricos) y GeoJSON
    geográfico (geometría con métricas operativas).}
    \label{fig:apimetro-salidas-tipos}
    \fuentefigura{Elaboración propia.}
\end{figure}

\texttt{Apimetro} ofrece dos tipos de respuesta según lo que el investigador o
la herramienta necesite, como ilustra la Figura~
\ref{fig:apimetro-salidas-tipos} . La primera —el JSON descriptivo— devuelve
únicamente atributos alfanuméricos: nombre comercial de cada línea, sistema,
color institucional y año de inauguración \citep{bray2017json}. No contiene
geometría. Es la «ficha técnica» de la red, adecuada para poblar tablas
comparativas o configurar la simbología de un visor sin necesidad de transferir
coordenadas.

El campo \texttt{clasificacion} —que distingue infraestructura existente, futura
o eliminada según los datos de la \textsc{adip} \citep{adip_gtfs_2024}— es el
insumo que habilita a \textit{VFTModel} para comparar la red actual con
escenarios de expansión \citep{cabrera2026VFTModel}. Los ejemplos completos de
ambas salidas se presentan en el Anexo~\ref{anx:apimetro-salidas}.

